% 【研读实验室】所属：第4章（由 chapters/revised/04-sources.tex \input 引入）
\minitopic{研读实验室：用同一组案例检验知觉理论}\index{意向论}\index{析取论}\index{认知渗透}\index{知觉学习} 知觉理论至少回答三个不同问题：经验是什么，经验如何把主体连接到对象，经验如何提供理由。直接实在论、意向论、感觉材料论与析取论对这些问题的组合不同。仅问“我们是否直接看见世界”容易把形而上学、心理学和认识论混在一起。 统一案例组可以减少术语争执。A是正常白昼看见红苹果；B是暖色灯下把白苹果看成红；C是无对象的红苹果幻觉；D是看见屏幕生成的红苹果图像；E是通过光谱仪读数判断苹果反射峰值。同一理论要说明五种经验的共同点、差异与理由地位，而不能只处理A。 感觉材料论可用共同内在对象解释A至C的现象相似，却面临如何从感觉材料到外部苹果的帷幕问题。意向论说经验以某种表征内容呈现世界，即使对象不存在也可有相同内容；它需要说明表征准确性与知觉理由的联系。析取论把A的成功对象接触与C的幻觉视为根本不同种类，能让事实本身进入理由，却要解释主体为何无法内省分辨。 直接实在论并不承诺永不错误。主体可直接面对真实对象，却因光线、注意或分类失误而判断错。把“直接”理解成“不可错”，会用认识结论替代关系说明。相反，承认加工与神经机制也不自动推出只知内部表象；中介因果机制与认识对象不是同一划分。

\minitopic{错觉、幻觉与误认}错觉中通常有对象，但某属性被错误呈现；幻觉中相关对象不存在；误认中经验可能准确，主体却把对象归入错误类别。远处将陌生人看成朋友多是误认：确实看见一个人，轮廓也许无误，错误发生在身份判断。不同错误对确证的意义不同。 若理论从“幻觉主观上不可区分”推出“真知觉与幻觉提供完全相同理由”，需要一项桥梁：主体可得理由只由内在现象决定。内在主义者可能接受；析取论者会说，真知觉中主体拥有由事实构成的更好理由，只是不一定知道自己拥有。这里分歧不是是否承认主观相似，而是相似是否穷尽认识地位。 错觉也未必取消全部知识。水中直棒看似弯曲，主体知道折射条件后仍可通过经验知道棒在那里，并合理地不相信它弯。经验提供多层信息：对象存在、表面呈现和背景条件。成熟感知者能用背景知识校正属性判断，不必把经验整体丢弃。 误认揭示概念与社会学习。证人准确看见面孔，却因照片阵列提示把它归为某嫌疑人；医生准确看到阴影，却因训练不足误诊。感官输入可以正常，推理与分类仍失败。说“眼见为实”或“眼睛欺骗”都太粗。

\minitopic{认知渗透与证据污染}若愿望、偏见或背景信念改变知觉经验本身，就发生认知渗透；若经验不变，只影响注意、记忆或判断，则是后知觉效应。二者在实验上难区分。主体更快发现自己期待的对象，可能因为视觉表征改变，也可能因为决策阈值降低。 认识论关心恶性渗透。若主体因相信某群体危险而把模糊动作看成威胁，再把“我看见威胁”当作原信念的新证据，形成自我支持循环。经验似乎提供独立确认，实际上由待证信念塑造。并非所有渗透都恶性：专家知识使放射科医生在影像中看出结构，可能增加真实信息敏感性。 区分标准包括真理联系、可校正性与来源独立性。训练若让经验对病灶变化更敏感、错误可由反馈纠正，具有认识增益；偏见若只提高某类判断阈值且抵抗反证，则具有污染风险。关键不是经验是否受上层影响，而是影响是否使系统更好响应事实。 制度可以减少污染。证人列队双盲、图像随机化、实验分析预注册、医生先独立判读再讨论，都在切断期望到判断的某些路径。它们不假设个体毫无偏见，而是让错误机制不易自我隐藏。

\minitopic{无概念经验与学习}经验能否提供主体尚无概念的内容？婴儿和动物能辨别颜色、数量与面孔，却未掌握成人词汇，这支持某种无概念内容。反对者说，构成理由的经验必须能进入判断空间；纯感觉差异若不能作为“如此”的理由，至多是原因。 争论需要区分内容拥有与内容报告。主体不能说出“洋红”，仍能可靠把样本分类；这显示缺少语言标签不等于没有任何表征能力。但若主体完全不能跨情形重认差异，称经验具有可作为理由的特定内容就较困难。概念可被理解为语言词、推理能力或辨别能力，不同定义会改变结论。 知觉学习说明概念与经验可能共同发展。品酒者、鸟类观察者和医学专家经过反馈后看到新差异。学习既可能改变早期感知敏感性，也可能只改变注意与分类。无论机制在哪一层，知识不是把固定感觉输入交给抽象理性；训练会塑造什么成为显著证据。 专家观看也需防止封闭。高熟练者可能更准确，却因权威地位而少受复核；共享训练可能使整个群体有共同盲点。将新手与专家表现、不同训练传统和盲法结果比较，才能证明专家经验确实提高事实响应。

\minitopic{比较视角：蔽、名与校正观看}《荀子》有以“蔽”讨论认识片面性的篇章：主体可能执着一端，以局部遮蔽整体；心的功能不是生成一个脱离世界的纯内在图像，而要统摄感官材料、名物区分与公共辨正\parencite{xunzi2014}。这与当代注意偏差和认知渗透可形成问题对话。 但“蔽”不等于实验心理学意义的视觉错觉。它常涉及价值执着、概念片面与论辩立场，范围比感官错误更广。若直接用现代神经术语解释，会丢失文本的修养与政治维度。比较时应先保持层次，再问共同问题：为何主体会把局部呈现误当全部，何种训练能让其响应被忽略的方面？ 当代知觉论提供精细区分：对象缺失、属性错呈、身份误认、注意选择和判断偏见。荀子式视角则提醒，校正观看不仅是个人感官校准，也依赖语言规范、师友反馈与去除固执。结合二者可形成多层纠错模型，而无需宣称它们拥有同一理论。

\minitopic{综合案例：卫星图像中的山火}分析员在增强后的卫星图像中看见疑似烟羽。原始光谱信号微弱，算法依据历史火灾模式提高对比；当地当日确有山火。分析员是否通过图像知道山火发生，取决于算法是否只是显化信号，还是把模型预测以像素形式重新写入图像。 若增强层把模型推断标成观测颜色，使用者可能把推论当直接视觉证据。真实火灾使结论为真，却不保证这一图像路径不含循环：历史模型预测火灾，增强制造烟羽外观，外观又被当作模型的独立验证。界面应区分原始测量、处理步骤与推断覆盖层。 独立证据可来自热红外、地面报告和另一颗传感器，但“独立”需检查共享数据与共同天气模型。分析员的专家能力包括识别处理伪影、比较波段、请求原始数据，而非只在屏幕上看得更快。人—仪器系统的知识要由整条校准链支持。 直接实在论可说仪器扩展了与山火的因果接触；意向论分析图像经验的内容与准确条件；析取论强调成功观测与无火伪影的根本差异；证据论则问屏幕呈现给予何种可被击败的理由。案例不必裁定唯一理论，却能显示每种理论实际解释什么。

\begin{tcolorbox}[breakable,colback=white,colframe=blue!30!black,title={本章研读任务},fonttitle=\bfseries]
为一个知觉判断制作五版：正常知觉、属性错觉、无对象幻觉、对象正确但身份误认、技术增强后的观测。逐版区分对象关系、经验内容、信念基础与击败者，并比较意向论、析取论和直接实在论的说明重点。
\end{tcolorbox}
